\chapter*{Nota de método}
\addcontentsline{toc}{chapter}{Nota de método}

Este caderno ficha fontes primárias, não literatura. Reúne o que os atores do
debate sobre a Caixa de Conversão disseram, em que data e sobre qual objeto, com
a citação direta e a localização na fonte. Serve ao mapeamento descritivo da
dissertação (debates, agentes, posições e marcos), e não à estimativa
quantitativa da posição dos jornais, que depende de instrumento ainda não
escolhido. Quando um jornal aparece aqui com voz própria, o registro é
inventário de voz com citação, não medida de posição editorial.

\bloco{Fontes}

Usei duas bases. A primeira são os dois volumes dos \textit{Documentos
Parlamentares} dedicados à Caixa de Conversão, impressos em Paris em 1914: o
volume de 1906, com a tramitação da lei de criação \cite{brasil1906caixa}, e o
de 1910, com a tramitação da lei que elevou a taxa a 16 dinheiros e o limite dos
depósitos \cite{brasil1910taxa16}. Li-os na camada de texto corrigido do projeto
e cito pela página impressa. A segunda base é o texto embutido (OCR da
Biblioteca Nacional) de quatro diários da Hemeroteca Digital Brasileira:
\textit{O Paiz}, \textit{Correio da Manhã} e \textit{Gazeta de Notícias}, do Rio
de Janeiro, e \textit{Correio Paulistano}, de São Paulo. Não há volume
parlamentar digitalizado para 1907 a 1909 nem para 1911 a 1914; nesses anos as
fichas dependem da imprensa.

\bloco{Convenções de citação}

As citações diretas preservam a grafia de época. No corpo do texto a ortografia
é a atual e os nomes seguem a forma corrente (Serzedelo Corrêa, Rui Barbosa).
Elipses e interpolações minhas vêm entre colchetes.

Toda citação direta foi conferida mecanicamente contra o texto da fonte, pela
regra do projeto: normalizada (sem acento nem pontuação), a citação tem de ser
subcadeia literal da página indicada. A citação que não casa literalmente, mas
alcança similaridade de janela de pelo menos 0,85, fica marcada com adaga
(\textsuperscript{\dag}): é trecho lido sobre OCR ruidoso ou cortado por quebra
de coluna, e pede conferência na imagem antes de entrar no texto da dissertação.
Citação abaixo desse limiar é rejeitada e não entra. O resultado de cada
citação está em \texttt{relatorio-conferencia-citacoes.md}, ao lado deste
arquivo.

Conferi na imagem da página, além da conferência mecânica, os pontos que
sustentam afirmações de autoria ou de data: o voto nominal de Alcindo Guanabara
em 14 de setembro de 1906, a atribuição da fala de Serzedelo Corrêa na
\textit{Gazeta de Notícias} de 12 de dezembro de 1914, os dois artigos sem
assinatura de \textit{O Paiz} de 1906, a lista de jornais na fala de Pinheiro
Machado em 1913 e o cabeçalho das edições cuja data o OCR não trazia. Esses
casos estão indicados nas fichas.

\bloco{Voz}

Relato de sessão parlamentar publicado em jornal é fala de terceiro: registra o
que o orador disse segundo o resumo do jornal, não a posição do jornal. O mesmo
vale para entrevista, carta e seção paga. Cada ficha distingue a fala direta
(anais, entrevista em primeira pessoa, carta), o resumo de jornal e o texto da
própria redação.

\bloco{Organização}

As fichas seguem os marcos do período: criação (1906), operação (1907 a 1909),
elevação da taxa e do limite (1910), crise (1913) e emissão com suspensão do
troco (1914). Duas fichas tratam de 1906, uma para os vencidos e outra para os
defensores, e uma terceira reúne o que se passou fora do plenário naquele ano.
A oitava ficha registra a voz própria dos jornais. A síntese final lista os
achados e o que continua em aberto.
